\problema{11}{Valor propio en 2D}{4}

\begin{enunciado}
	Para la malla del Problema 6 sea $\mat{A}\vect{\varphi} = \frac1k\mat{F}\vect{\varphi}$ con $\mat{A}$
	del Problema 6 y $\mat{F} = 3\mat{I}$.
	\begin{enumerate}[label=(\alph*)]
		\item Obtenga los valores propios de $\mat{A}$, los de $k$, $\keff$ y $\rho_D$.
		\item Realice una iteración de potencias desde $\vect{\varphi}^{(0)} = (1, 0, 0, 0)^T$, $k^{(0)} = 1$.
		\item Afecta a la convergencia que el segundo valor propio tenga multiplicidad 2? Explique.
	\end{enumerate}
\end{enunciado}

Del Problema 6
\[
\mat{A} =
\begin{pmatrix}
	4 & -1 & -1 &  0\\
	-1 &  4 &  0 & -1\\
	-1 &  0 &  4 & -1\\
	0 & -1 & -1 &  4
\end{pmatrix}
= \begin{pmatrix} \mat{T} & -\mat{I}\\ -\mat{I} & \mat{T} \end{pmatrix},
\qquad \mat{T} = \begin{pmatrix} 4 & -1\\ -1 & 4 \end{pmatrix},
\qquad \mat{F} = 3\mat{I} .
\]

\inciso{a}{Valores propios de \texorpdfstring{$\mat{A}$}{A}, de \texorpdfstring{$k$}{k}, \texorpdfstring{$\keff$}{kef} y \texorpdfstring{$\rho_D$}{rhoD}}


\paso{1}{Reducir el problema}
Se sustituye $\mat{F} = 3\mat{I}$:
\[
\mat{A}\vect{\varphi} = \frac1k\,(3\mat{I})\,\vect{\varphi} = \frac3k\,(\mat{I}\vect{\varphi}) .
\]
Como $\mat{I}\vect{\varphi} = \vect{\varphi}$ (a identidad no cambia al vector
\[
3\mat{I}\vect{\varphi} =
\begin{pmatrix} 3 & 0 & 0 & 0\\ 0 & 3 & 0 & 0\\ 0 & 0 & 3 & 0\\ 0 & 0 & 0 & 3 \end{pmatrix}
\begin{pmatrix} \varphi_1\\ \varphi_2\\ \varphi_3\\ \varphi_4 \end{pmatrix}
= \begin{pmatrix} 3\varphi_1\\ 3\varphi_2\\ 3\varphi_3\\ 3\varphi_4 \end{pmatrix} = 3\vect{\varphi}
\qquad\Longrightarrow\qquad
\mat{A}\vect{\varphi} = \frac3k\,\vect{\varphi} .
\]
Esta es la definición de valor propio de $\mat{A}$: si $\mat{A}\vect{\varphi} = \mu\,\vect{\varphi}$, entonces
\[
\mu = \frac3k \qquad\Longrightarrow\qquad k = \frac{3}{\mu},
\]
y los vectores propios son los mismos que los de $\mat{A}$.

Para el método de potencias conviene que $k$ quede como valor propio
\[
\vect{\varphi} = \frac1k\,\mat{A}^{-1}\mat{F}\vect{\varphi}
\qquad\Longrightarrow\qquad
\underbrace{\mat{A}^{-1}\mat{F}}_{\mat{M}}\,\vect{\varphi} = k\,\vect{\varphi},
\qquad \mat{M} = 3\mat{A}^{-1} .
\]
Si $\mat{A}\vect{\varphi} = \mu\vect{\varphi}$, al multiplicar por $\mat{A}^{-1}$ queda $\vect{\varphi} = \mu\,\mat{A}^{-1}\vect{\varphi}$, o sea
$\mat{A}^{-1}\vect{\varphi} = \frac1\mu\vect{\varphi}$, y entonces
\[
\mat{M}\vect{\varphi} = 3\mat{A}^{-1}\vect{\varphi} = \frac3\mu\,\vect{\varphi} .
\]
Los valores propios de $\mat{M}$ son $k = 3/\mu$ por lo que hay que encontrar los valores propios de $\mat{A}$.

\paso{2}{Condición para los valores propios}
$\mat{A}\vect{\varphi} = \mu\vect{\varphi}$ con $\vect{\varphi} \neq \vect{0}$. Pasando todo a un lado y usando
$\mu\vect{\varphi} = \mu\mat{I}\vect{\varphi}$:
\[
(\mat{A} - \mu\mat{I})\,\vect{\varphi} = \vect{0} .
\]
Si $\det(\mat{A} - \mu\mat{I}) \neq 0$ la única solución es $\vect{\varphi} = \vect{0}$ (no sirve). Para tener
$\vect{\varphi} \neq \vect{0}$ se necesita
\[
\det(\mat{A} - \mu\mat{I}) = 0 .
\]

\paso{3}{Escribir $\mat{A} - \mu\mat{I}$ por bloques}
Restar $\mu\mat{I}$ solo cambia la diagonal. Con $a = 4 - \mu$:
\[
\mat{A} - \mu\mat{I} =
\left(\begin{array}{cc|cc}
	a & -1 & -1 &  0\\
	-1 &  a &  0 & -1\\ \hline
	-1 &  0 &  a & -1\\
	0 & -1 & -1 &  a
\end{array}\right)
= \begin{pmatrix} \mat{T} - \mu\mat{I} & -\mat{I}\\ -\mat{I} & \mat{T} - \mu\mat{I} \end{pmatrix},
\qquad
\mat{T} - \mu\mat{I} = \begin{pmatrix} a & -1\\ -1 & a \end{pmatrix}.
\]

\paso{4}{Determinante por bloques}

\[
\det(\mat{A} - \mu\mat{I})
= \det\big[(\mat{T} - \mu\mat{I})(\mat{T} - \mu\mat{I}) - (-\mat{I})(-\mat{I})\big]
= \det\big[(\mat{T} - \mu\mat{I})^2 - \mat{I}\big].
\]
Se eleva al cuadrado (renglón por columna):
\[
(\mat{T} - \mu\mat{I})^2 =
\begin{pmatrix} a & -1\\ -1 & a \end{pmatrix}
\begin{pmatrix} a & -1\\ -1 & a \end{pmatrix}
= \begin{pmatrix} a^2 + 1 & -a - a\\ -a - a & 1 + a^2 \end{pmatrix}
= \begin{pmatrix} a^2 + 1 & -2a\\ -2a & a^2 + 1 \end{pmatrix}.
\]
Se resta $I$
\[
(\mat{T} - \mu\mat{I})^2 - \mat{I} = \begin{pmatrix} a^2 & -2a\\ -2a & a^2 \end{pmatrix}.
\]
Su determinante:
\[
\det(\mat{A} - \mu\mat{I}) = a^2\cdot a^2 - (-2a)(-2a) = a^4 - 4a^2 = a^2\left(a^2 - 4\right) = a^2\,(a - 2)(a + 2).
\]

\paso{5}{Igualar a cero}
\[
a^2\,(a - 2)(a + 2) = 0, \qquad a = 4 - \mu \ \Longrightarrow\ \mu = 4 - a .
\]
\begin{center}
	\begin{tabular}{|c|c|c|}
		\hline
		Factor & $a$ & $\mu = 4 - a$\\
		\hline
		$a^2 = 0$ & $0$ & $4$ \\
		$a - 2 = 0$ & $2$ & $2$\\
		$a + 2 = 0$ & $-2$ & $6$\\
		\hline
	\end{tabular}
\end{center}
\resultado{\mu_1 = 2, \qquad \mu_2 = \mu_3 = 4, \qquad \mu_4 = 6}

\paso{6}{Vectores propios}
Por la simetria
\begin{align*}
	\mat{A}\begin{pmatrix} 1\\ 1\\ 1\\ 1 \end{pmatrix}
	&= \begin{pmatrix} 4 - 1 - 1 + 0\\ -1 + 4 + 0 - 1\\ -1 + 0 + 4 - 1\\ 0 - 1 - 1 + 4 \end{pmatrix}
	= \begin{pmatrix} 2\\ 2\\ 2\\ 2 \end{pmatrix} = 2\begin{pmatrix} 1\\ 1\\ 1\\ 1 \end{pmatrix}\\[6pt]
	\mat{A}\begin{pmatrix} 1\\ -1\\ 1\\ -1 \end{pmatrix}
	&= \begin{pmatrix} 4 + 1 - 1 + 0\\ -1 - 4 + 0 + 1\\ -1 + 0 + 4 + 1\\ 0 + 1 - 1 - 4 \end{pmatrix}
	= \begin{pmatrix} 4\\ -4\\ 4\\ -4 \end{pmatrix} = 4\begin{pmatrix} 1\\ -1\\ 1\\ -1 \end{pmatrix}\\[6pt]
	\mat{A}\begin{pmatrix} 1\\ 1\\ -1\\ -1 \end{pmatrix}
	&= \begin{pmatrix} 4 - 1 + 1 + 0\\ -1 + 4 + 0 + 1\\ -1 + 0 - 4 + 1\\ 0 - 1 + 1 - 4 \end{pmatrix}
	= \begin{pmatrix} 4\\ 4\\ -4\\ -4 \end{pmatrix} = 4\begin{pmatrix} 1\\ 1\\ -1\\ -1 \end{pmatrix}\\[6pt]
	\mat{A}\begin{pmatrix} 1\\ -1\\ -1\\ 1 \end{pmatrix}
	&= \begin{pmatrix} 4 + 1 + 1 + 0\\ -1 - 4 + 0 - 1\\ -1 + 0 - 4 - 1\\ 0 + 1 + 1 + 4 \end{pmatrix}
	= \begin{pmatrix} 6\\ -6\\ -6\\ 6 \end{pmatrix} = 6\begin{pmatrix} 1\\ -1\\ -1\\ 1 \end{pmatrix}
\end{align*}

\paso{7}{Valores de $k$ y $\keff$}

Con $k_i = 3/\mu_i$:
\begin{center}
	\begin{tabular}{|c|c|c|c|}
		\hline
		Modo & Vector propio & $\mu_i$ & $k_i = 3/\mu_i$\\
		\hline
		1 & $(1, 1, 1, 1)^T$   & 2 & $3/2 $\\
		2 & $(1, -1, 1, -1)^T$ & 4 & $3/4 $\\
		3 & $(1, 1, -1, -1)^T$ & 4 & $3/4 $\\
		4 & $(1, -1, -1, 1)^T$ & 6 & $3/6 = 1/2$\\
		\hline
	\end{tabular}
\end{center}
Aqui es curioso que el orden se invierte, el valor propio mas chico de $\mat{A}$ da el $k$ más grande

\resultado{\keff = k_1 = \frac{3}{\mu_{\min}} = \frac32 }
Corresponde al modo con todos los nodos positivos (Perron--Frobenius) $\keff$ es real, positivo y simple

\paso{8}{Razón de dominancia}
\[
\rho_D = \frac{k_2}{k_1} = \frac{\frac34}{\frac32} = \frac{\mu_1}{\mu_2} = \frac24 = \frac12
\]

% =====================================================================
\inciso{b}{Una iteración de potencias}
% =====================================================================

\[
\vect{q}^{(n)} = \frac{1}{k^{(n-1)}}\,\mat{F}\vect{\varphi}^{(n-1)},
\qquad
\mat{A}\vect{\varphi}^{(n)} = \vect{q}^{(n)},
\qquad
k^{(n)} = k^{(n-1)}\,\frac{\vect{1}^T\mat{F}\vect{\varphi}^{(n)}}{\vect{1}^T\mat{F}\vect{\varphi}^{(n-1)}},
\]
se divide entre su maximo

Datos: $\vect{\varphi}^{(0)} = (1, 0, 0, 0)^T$, $k^{(0)} = 1$, $\mat{F} = 3\mat{I}$.

\paso{1}{Fuente}
\[
\vect{q}^{(1)} = \frac{1}{k^{(0)}}\,\mat{F}\vect{\varphi}^{(0)} = \frac11\cdot 3\begin{pmatrix} 1\\ 0\\ 0\\ 0 \end{pmatrix}
= \begin{pmatrix} 3\\ 0\\ 0\\ 0 \end{pmatrix}.
\]

Resolver $\mat{A}\vect{\varphi}^{(1)} = \vect{q}^{(1)}$ con el factor $\mat{G}$ del Problema 6
\[
\mat{A}\vect{\varphi}^{(1)} = \vect{q}^{(1)}
\;\Longrightarrow\;
\mat{G}\underbrace{\mat{G}^T\vect{\varphi}^{(1)}}_{\vect{y}} = \vect{q}^{(1)}
\;\Longrightarrow\;
\begin{cases}
	\mat{G}\vect{y} = \vect{q}^{(1)} & \text{(sustitución hacia adelante)}\\
	\mat{G}^T\vect{\varphi}^{(1)} = \vect{y} & \text{(sustitución hacia atrás)}
\end{cases}
\]
con
\[
g_{11} = 2,\quad g_{21} = g_{31} = -\tfrac12,\quad g_{41} = 0,\quad g_{22} = \tfrac{\sqrt{15}}{2},\quad
g_{32} = -\tfrac{1}{2\sqrt{15}},\quad g_{42} = -\tfrac{2}{\sqrt{15}},
\]
\[
g_{33} = \sqrt{\tfrac{56}{15}},\quad g_{43} = -\tfrac{8}{\sqrt{210}},\quad g_{44} = \sqrt{\tfrac{24}{7}} .
\]

\textbf{Hacia adelante.} El sistema $\mat{G}\vect{y} = \vect{q}^{(1)}$ es
\[
\begin{pmatrix}
	g_{11} & 0 & 0 & 0\\
	g_{21} & g_{22} & 0 & 0\\
	g_{31} & g_{32} & g_{33} & 0\\
	g_{41} & g_{42} & g_{43} & g_{44}
\end{pmatrix}
\begin{pmatrix} y_1\\ y_2\\ y_3\\ y_4 \end{pmatrix}
=
\begin{pmatrix} 3\\ 0\\ 0\\ 0 \end{pmatrix}
\quad\Longrightarrow\quad
\begin{aligned}
	g_{11}y_1 &= 3\\
	g_{21}y_1 + g_{22}y_2 &= 0\\
	g_{31}y_1 + g_{32}y_2 + g_{33}y_3 &= 0\\
	g_{41}y_1 + g_{42}y_2 + g_{43}y_3 + g_{44}y_4 &= 0
\end{aligned}
\]
Se despeja de arriba hacia abajo:
\begin{align*}
	y_1 &= \frac{3}{g_{11}} = \frac32\\[8pt]
	y_2 &= \frac{0 - g_{21}y_1}{g_{22}}
	= \frac{0 - \left(-\tfrac12\right)\left(\tfrac32\right)}{\tfrac{\sqrt{15}}{2}}
	= \frac{\tfrac34}{\tfrac{\sqrt{15}}{2}} = \frac{3}{2\sqrt{15}} = \frac{\sqrt{15}}{10}\\[8pt]
	y_3 &= \frac{0 - g_{31}y_1 - g_{32}y_2}{g_{33}}
	= \frac{0 - \left(-\tfrac12\right)\left(\tfrac32\right) - \left(-\tfrac{1}{2\sqrt{15}}\right)\left(\tfrac{\sqrt{15}}{10}\right)}{\sqrt{\tfrac{56}{15}}}
	= \frac{\tfrac34 + \tfrac{1}{20}}{\sqrt{\tfrac{56}{15}}}
	= \frac{\tfrac45}{\sqrt{\tfrac{56}{15}}} = \frac{\sqrt{210}}{35}\\[8pt]
	y_4 &= \frac{0 - g_{41}y_1 - g_{42}y_2 - g_{43}y_3}{g_{44}}
	= \frac{0 - 0 - \left(-\tfrac{2}{\sqrt{15}}\right)\left(\tfrac{\sqrt{15}}{10}\right) - \left(-\tfrac{8}{\sqrt{210}}\right)\left(\tfrac{\sqrt{210}}{35}\right)}{\sqrt{\tfrac{24}{7}}}
	= \frac{\tfrac15 + \tfrac{8}{35}}{\sqrt{\tfrac{24}{7}}}
	= \frac{\tfrac37}{\sqrt{\tfrac{24}{7}}} = \frac{\sqrt{42}}{28}
\end{align*}

\textbf{Hacia atrás.} El sistema $\mat{G}^T\vect{\varphi}^{(1)} = \vect{y}$ es
\[
\begin{pmatrix}
	g_{11} & g_{21} & g_{31} & g_{41}\\
	0 & g_{22} & g_{32} & g_{42}\\
	0 & 0 & g_{33} & g_{43}\\
	0 & 0 & 0 & g_{44}
\end{pmatrix}
\begin{pmatrix} \varphi_1^{(1)}\\ \varphi_2^{(1)}\\ \varphi_3^{(1)}\\ \varphi_4^{(1)} \end{pmatrix}
=
\begin{pmatrix} y_1\\ y_2\\ y_3\\ y_4 \end{pmatrix}
\quad\Longrightarrow\quad
\begin{aligned}
	g_{11}\varphi_1^{(1)} + g_{21}\varphi_2^{(1)} + g_{31}\varphi_3^{(1)} + g_{41}\varphi_4^{(1)} &= y_1\\
	g_{22}\varphi_2^{(1)} + g_{32}\varphi_3^{(1)} + g_{42}\varphi_4^{(1)} &= y_2\\
	g_{33}\varphi_3^{(1)} + g_{43}\varphi_4^{(1)} &= y_3\\
	g_{44}\varphi_4^{(1)} &= y_4
\end{aligned}
\]
Se despeja de abajo hacia arriba
\begin{align*}
	\varphi_4^{(1)} &= \frac{y_4}{g_{44}} = \frac{\tfrac{\sqrt{42}}{28}}{\sqrt{\tfrac{24}{7}}}
	= \frac{1}{28}\sqrt{\frac{42\cdot 7}{24}} = \frac{1}{28}\sqrt{\frac{49}{4}} = \frac{1}{28}\cdot\frac72 = \frac18\\[8pt]
	\varphi_3^{(1)} &= \frac{y_3 - g_{43}\varphi_4^{(1)}}{g_{33}}
	= \frac{\tfrac{\sqrt{210}}{35} - \left(-\tfrac{8}{\sqrt{210}}\right)\left(\tfrac18\right)}{\sqrt{\tfrac{56}{15}}}
	= \frac{\tfrac{\sqrt{210}}{35} + \tfrac{\sqrt{210}}{210}}{\sqrt{\tfrac{56}{15}}}
	= \frac{\tfrac{\sqrt{210}}{30}}{\sqrt{\tfrac{56}{15}}} = \frac{1}{30}\sqrt{\frac{210\cdot15}{56}} = \frac{1}{30}\cdot\frac{15}{2} = \frac14\\[8pt]
	\varphi_2^{(1)} &= \frac{y_2 - g_{32}\varphi_3^{(1)} - g_{42}\varphi_4^{(1)}}{g_{22}}
	= \frac{\tfrac{\sqrt{15}}{10} - \left(-\tfrac{1}{2\sqrt{15}}\right)\left(\tfrac14\right) - \left(-\tfrac{2}{\sqrt{15}}\right)\left(\tfrac18\right)}{\tfrac{\sqrt{15}}{2}}
	= \frac{\tfrac{12}{8\sqrt{15}} + \tfrac{1}{8\sqrt{15}} + \tfrac{2}{8\sqrt{15}}}{\tfrac{\sqrt{15}}{2}}
	= \frac{\tfrac{15}{8\sqrt{15}}}{\tfrac{\sqrt{15}}{2}} = \frac14\\[8pt]
	\varphi_1^{(1)} &= \frac{y_1 - g_{21}\varphi_2^{(1)} - g_{31}\varphi_3^{(1)} - g_{41}\varphi_4^{(1)}}{g_{11}}
	= \frac{\tfrac32 - \left(-\tfrac12\right)\left(\tfrac14\right) - \left(-\tfrac12\right)\left(\tfrac14\right) - 0}{2}
	= \frac{\tfrac32 + \tfrac18 + \tfrac18}{2} = \frac{\tfrac74}{2} = \frac78
\end{align*}
(En $\varphi_2^{(1)}$ se usó $\tfrac{\sqrt{15}}{10} = \tfrac{15}{10\sqrt{15}} = \tfrac{12}{8\sqrt{15}}$.)
\[
\vect{\varphi}^{(1)} = \left(\tfrac78,\ \tfrac14,\ \tfrac14,\ \tfrac18\right)^T
\]

\paso{3}{Actualizar $k$}
\begin{align*}
	\vect{1}^T\mat{F}\vect{\varphi}^{(0)} &= 3\,(1 + 0 + 0 + 0) = 3\\
	\vect{1}^T\mat{F}\vect{\varphi}^{(1)} &= 3\left(\tfrac78 + \tfrac14 + \tfrac14 + \tfrac18\right) = 3\left(\tfrac{12}{8}\right) = \tfrac92
\end{align*}
\[
k^{(1)} = k^{(0)}\,\frac{\vect{1}^T\mat{F}\vect{\varphi}^{(1)}}{\vect{1}^T\mat{F}\vect{\varphi}^{(0)}}
= 1\cdot\frac{9/2}{3} = \frac32 .
\]

Se divide entre el maximo, $\max_i \varphi_i^{(1)} = \tfrac78$:
\[
\frac{1}{7/8}\left(\tfrac78,\ \tfrac14,\ \tfrac14,\ \tfrac18\right)^T
= \left(\tfrac88,\ \tfrac{8}{28},\ \tfrac{8}{28},\ \tfrac{8}{56}\right)^T
= \left(1,\ \tfrac27,\ \tfrac27,\ \tfrac17\right)^T
\]
\resultado{\vect{\varphi}^{(1)} = \left(1,\ \tfrac27,\ \tfrac27,\ \tfrac17\right)^T \approx (1,\ 0.285714,\ 0.285714,\ 0.142857)^T}

\paso{5}{Criterio de convergencia}
\[
\frac{\left|k^{(1)} - k^{(0)}\right|}{k^{(1)}} = \frac{\left|\frac32 - 1\right|}{\frac32} = 0.333333 > 10^{-5},
\]
no ha convergido


\inciso{c}{Afecta la multiplicidad 2 a la convergencia?}
% =====================================================================
No. En cada iteración, comparada con el modo fundamental, la parte de cada uno de los otros modos se multiplica
por $k_i/k_1$